\hypertarget{tt__diagram_8c}{
\section{tt\_\-diagram.c File Reference}
\label{tt__diagram_8c}\index{tt_diagram.c@{tt\_\-diagram.c}}
}
plugin that creates a file-info \char`\"{}tooltip\char`\"{} 

\subsection*{Defines}
\begin{CompactItemize}
\item 
\hypertarget{tt__diagram_8c_a3}{
\#define \hyperlink{tt__diagram_8c_a3}{WITH\_\-THUMBNAILS}}
\label{tt__diagram_8c_a3}

\begin{CompactList}\small\item\em define for funky thumbnails \item\end{CompactList}\item 
\#define \hyperlink{tt__diagram_8c_a4}{COL}~255
\begin{CompactList}\small\item\em valid colours are 0 and 255. \item\end{CompactList}\item 
\hypertarget{tt__diagram_8c_a5}{
\#define \hyperlink{tt__diagram_8c_a5}{LEN}~24}
\label{tt__diagram_8c_a5}

\begin{CompactList}\small\item\em string-representation of a value can not get longer than this \item\end{CompactList}\item 
\#define \hyperlink{tt__diagram_8c_a6}{LINES\_\-FOR\_\-ALL}~(RND(3)$<$1)
\begin{CompactList}\small\item\em draw lines for each text-line (1), not just for each text-block (0) instead, a function may be given to randomly select text-lines that will get a line-object. \item\end{CompactList}\item 
\hypertarget{tt__diagram_8c_a7}{
\#define \hyperlink{tt__diagram_8c_a7}{VAL\_\-ALIGN}}
\label{tt__diagram_8c_a7}

\begin{CompactList}\small\item\em define for all values to be aligned (same X-cordinate) \item\end{CompactList}\item 
\hypertarget{tt__diagram_8c_a8}{
\#define \hyperlink{tt__diagram_8c_a8}{BOTTOM\_\-ALIGN}}
\label{tt__diagram_8c_a8}

\begin{CompactList}\small\item\em define for key and value to share a baseline. \item\end{CompactList}\item 
\hypertarget{tt__diagram_8c_a9}{
\#define \hyperlink{tt__diagram_8c_a9}{KVS}~24}
\label{tt__diagram_8c_a9}

\begin{CompactList}\small\item\em padding between key and value \item\end{CompactList}\item 
\hypertarget{tt__diagram_8c_a10}{
\#define \hyperlink{tt__diagram_8c_a10}{BOX}}
\label{tt__diagram_8c_a10}

\begin{CompactList}\small\item\em draw frame around text-blocks? \item\end{CompactList}\item 
\hypertarget{tt__diagram_8c_a11}{
\#define \hyperlink{tt__diagram_8c_a11}{BOX\_\-TITLES}}
\label{tt__diagram_8c_a11}

\begin{CompactList}\small\item\em put group-name in frame? \item\end{CompactList}\item 
\hypertarget{tt__diagram_8c_a12}{
\#define \hyperlink{tt__diagram_8c_a12}{PAD}~8}
\label{tt__diagram_8c_a12}

\begin{CompactList}\small\item\em pad frame this much (pad$>$0) \item\end{CompactList}\item 
\#define \hyperlink{tt__diagram_8c_a14}{LINES\_\-FOR\_\-ALL}~0
\begin{CompactList}\small\item\em draw lines for each text-line (1), not just for each text-block (0) instead, a function may be given to randomly select text-lines that will get a line-object. \item\end{CompactList}\end{CompactItemize}
\subsection*{Functions}
\begin{CompactItemize}
\item 
static \hyperlink{struct__icon__el}{icon\_\-el} $\ast$ \hyperlink{tt__diagram_8c_a18}{frame\_\-new} (Evas $\ast$evas, \hyperlink{struct__icon}{icon} $\ast$\hyperlink{struct__icon}{icon}, \hyperlink{struct__icon__el}{icon\_\-el} $\ast$el, double x1, double y1, double x2, double y2, int state)
\begin{CompactList}\small\item\em create a frame around a group of info-items. \item\end{CompactList}\item 
static \hyperlink{struct__icon__el}{icon\_\-el} $\ast$ \hyperlink{tt__diagram_8c_a19}{line\_\-new} (Evas $\ast$evas, \hyperlink{struct__icon}{icon} $\ast$\hyperlink{struct__icon}{icon}, \hyperlink{struct__icon__el}{icon\_\-el} $\ast$el, double x1, double y1, double x2, double y2, int state)
\begin{CompactList}\small\item\em create a line connecting the file-icon with an info-item. \item\end{CompactList}\item 
static \hyperlink{struct__icon}{icon} $\ast$ \hyperlink{tt__diagram_8c_a20}{tt\_\-icon\_\-new} (const char $\ast$label)
\begin{CompactList}\small\item\em create an \char`\"{}icon\char`\"{} (a group of info-lines) \item\end{CompactList}\item 
static int \hyperlink{tt__diagram_8c_a21}{align\_\-left} (\hyperlink{struct__icon__list}{icon\_\-list} $\ast$il, Evas $\ast$evas, struct \hyperlink{struct__icon}{\_\-icon} $\ast$\hyperlink{struct__icon}{icon}, struct \hyperlink{struct__icon}{\_\-icon} $\ast$i, \hyperlink{struct__icon__el}{icon\_\-el} $\ast$el2, int x, int y)
\begin{CompactList}\small\item\em put group of info-items to the left of file-icon. \item\end{CompactList}\item 
static int \hyperlink{tt__diagram_8c_a22}{align\_\-right} (\hyperlink{struct__icon__list}{icon\_\-list} $\ast$il, Evas $\ast$evas, struct \hyperlink{struct__icon}{\_\-icon} $\ast$\hyperlink{struct__icon}{icon}, struct \hyperlink{struct__icon}{\_\-icon} $\ast$i, \hyperlink{struct__icon__el}{icon\_\-el} $\ast$el2, int x, int y)
\begin{CompactList}\small\item\em put group of info-items to the right of file-icon. \item\end{CompactList}\item 
int \hyperlink{tt__diagram_8c_a23}{align\_\-values} (struct \hyperlink{struct__icon}{\_\-icon} $\ast$\hyperlink{struct__icon}{icon}, int wk)
\begin{CompactList}\small\item\em align all values of a group (same x-ordinate) \item\end{CompactList}\item 
void \hyperlink{tt__diagram_8c_a24}{tt\_\-dst} (void $\ast$$\ast$data)
\begin{CompactList}\small\item\em destruct a tooltip. \item\end{CompactList}\item 
void $\ast$ \hyperlink{tt__diagram_8c_a25}{tt\_\-new} (void $\ast$data, int x, int y, \hyperlink{struct__file__info}{file\_\-info} $\ast$fi)
\begin{CompactList}\small\item\em create a tooltip. \item\end{CompactList}\end{CompactItemize}


\subsection{Detailed Description}
plugin that creates a file-info \char`\"{}tooltip\char`\"{} 

this plugin creates tooltips as shown in \href{http://evidence.sourceforge.net/screenshots/tooltip2.jpg}{\tt http://evidence.sourceforge.net/screenshots/tooltip2.jpg}

(if you are looking for a simple tooltip plugin to modify, you might prefer tt\_\-traditional and friends; this one is slightly more involved because of the code that tries to decide which boxes should go where.)

a lot of tunables exist as defines and should be turned into properties at some point. \#\#\#

\subsection{Define Documentation}
\hypertarget{tt__diagram_8c_a4}{
\index{tt_diagram.c@{tt\_\-diagram.c}!COL@{COL}}
\index{COL@{COL}!tt_diagram.c@{tt\_\-diagram.c}}
\subsubsection[COL]{\setlength{\rightskip}{0pt plus 5cm}\#define COL~255}}
\label{tt__diagram_8c_a4}


valid colours are 0 and 255. 

choose any other value, and you're on your own. \hypertarget{tt__diagram_8c_a14}{
\index{tt_diagram.c@{tt\_\-diagram.c}!LINES_FOR_ALL@{LINES\_\-FOR\_\-ALL}}
\index{LINES_FOR_ALL@{LINES\_\-FOR\_\-ALL}!tt_diagram.c@{tt\_\-diagram.c}}
\subsubsection[LINES\_\-FOR\_\-ALL]{\setlength{\rightskip}{0pt plus 5cm}\#define LINES\_\-FOR\_\-ALL~0}}
\label{tt__diagram_8c_a14}


draw lines for each text-line (1), not just for each text-block (0) instead, a function may be given to randomly select text-lines that will get a line-object. 

the first text-line of a text-block always will. \hypertarget{tt__diagram_8c_a6}{
\index{tt_diagram.c@{tt\_\-diagram.c}!LINES_FOR_ALL@{LINES\_\-FOR\_\-ALL}}
\index{LINES_FOR_ALL@{LINES\_\-FOR\_\-ALL}!tt_diagram.c@{tt\_\-diagram.c}}
\subsubsection[LINES\_\-FOR\_\-ALL]{\setlength{\rightskip}{0pt plus 5cm}\#define LINES\_\-FOR\_\-ALL~(RND(3)$<$1)}}
\label{tt__diagram_8c_a6}


draw lines for each text-line (1), not just for each text-block (0) instead, a function may be given to randomly select text-lines that will get a line-object. 

the first text-line of a text-block always will. 

\subsection{Function Documentation}
\hypertarget{tt__diagram_8c_a21}{
\index{tt_diagram.c@{tt\_\-diagram.c}!align_left@{align\_\-left}}
\index{align_left@{align\_\-left}!tt_diagram.c@{tt\_\-diagram.c}}
\subsubsection[align\_\-left]{\setlength{\rightskip}{0pt plus 5cm}static int align\_\-left (\hyperlink{struct__icon__list}{icon\_\-list} $\ast$ {\em il}, Evas $\ast$ {\em evas}, struct \hyperlink{struct__icon}{\_\-icon} $\ast$ {\em icon}, struct \hyperlink{struct__icon}{\_\-icon} $\ast$ {\em i}, \hyperlink{struct__icon__el}{icon\_\-el} $\ast$ {\em el2}, int {\em x}, int {\em y})\hspace{0.3cm}{\tt  \mbox{[}inline, static\mbox{]}}}}
\label{tt__diagram_8c_a21}


put group of info-items to the left of file-icon. 

\begin{Desc}
\item[Returns:]y-coordinate for next box on this side \end{Desc}
\begin{Desc}
\item[Parameters:]
\begin{description}
\item[{\em il}]- the icon-list \item[{\em evas}]- the canvas \item[{\em icon}]- the group of info-items \item[{\em i}]- the file-icon \item[{\em el2}]- the current icon-element \item[{\em x}]- x-offset for all boxes on this side \item[{\em y}]- y-offset for current box on this side\end{description}
\end{Desc}
\hypertarget{tt__diagram_8c_a22}{
\index{tt_diagram.c@{tt\_\-diagram.c}!align_right@{align\_\-right}}
\index{align_right@{align\_\-right}!tt_diagram.c@{tt\_\-diagram.c}}
\subsubsection[align\_\-right]{\setlength{\rightskip}{0pt plus 5cm}static int align\_\-right (\hyperlink{struct__icon__list}{icon\_\-list} $\ast$ {\em il}, Evas $\ast$ {\em evas}, struct \hyperlink{struct__icon}{\_\-icon} $\ast$ {\em icon}, struct \hyperlink{struct__icon}{\_\-icon} $\ast$ {\em i}, \hyperlink{struct__icon__el}{icon\_\-el} $\ast$ {\em el2}, int {\em x}, int {\em y})\hspace{0.3cm}{\tt  \mbox{[}inline, static\mbox{]}}}}
\label{tt__diagram_8c_a22}


put group of info-items to the right of file-icon. 

\begin{Desc}
\item[Returns:]y-coordinate for next box on this side \end{Desc}
\begin{Desc}
\item[Parameters:]
\begin{description}
\item[{\em il}]- the icon-list \item[{\em evas}]- the canvas \item[{\em icon}]- the group of info-items \item[{\em i}]- the file-icon \item[{\em el2}]- the current icon-element \item[{\em x}]- x-offset for all boxes on this side \item[{\em y}]- y-offset for current box on this side\end{description}
\end{Desc}
\hypertarget{tt__diagram_8c_a23}{
\index{tt_diagram.c@{tt\_\-diagram.c}!align_values@{align\_\-values}}
\index{align_values@{align\_\-values}!tt_diagram.c@{tt\_\-diagram.c}}
\subsubsection[align\_\-values]{\setlength{\rightskip}{0pt plus 5cm}int align\_\-values (struct \hyperlink{struct__icon}{\_\-icon} $\ast$ {\em icon}, int {\em wk})}}
\label{tt__diagram_8c_a23}


align all values of a group (same x-ordinate) 

\begin{Desc}
\item[Returns:]nw - new width of the group \end{Desc}
\begin{Desc}
\item[Parameters:]
\begin{description}
\item[{\em icon}]- the group of info-items \item[{\em wk}]- widest key in group\end{description}
\end{Desc}
\hypertarget{tt__diagram_8c_a18}{
\index{tt_diagram.c@{tt\_\-diagram.c}!frame_new@{frame\_\-new}}
\index{frame_new@{frame\_\-new}!tt_diagram.c@{tt\_\-diagram.c}}
\subsubsection[frame\_\-new]{\setlength{\rightskip}{0pt plus 5cm}static \hyperlink{struct__icon__el}{icon\_\-el}$\ast$ frame\_\-new (Evas $\ast$ {\em evas}, \hyperlink{struct__icon}{icon} $\ast$ {\em icon}, \hyperlink{struct__icon__el}{icon\_\-el} $\ast$ {\em el}, double {\em x1}, double {\em y1}, double {\em x2}, double {\em y2}, int {\em state})\hspace{0.3cm}{\tt  \mbox{[}static\mbox{]}}}}
\label{tt__diagram_8c_a18}


create a frame around a group of info-items. 

\begin{Desc}
\item[Returns:]a linked list of icon-elements. \end{Desc}
\begin{Desc}
\item[Parameters:]
\begin{description}
\item[{\em evas}]- the evas \item[{\em icon}]- the tooltip to append the elements to (needed by ICON\_\-EL\_\-$\ast$()) \item[{\em el}]- the current icon-element \item[{\em x1}]- left boundary \item[{\em y1}]- upper boundary \item[{\em x2}]- right boundary \item[{\em y2}]- lower boundary \item[{\em state}]- state mask for the icon-elements to create\end{description}
\end{Desc}
\hypertarget{tt__diagram_8c_a19}{
\index{tt_diagram.c@{tt\_\-diagram.c}!line_new@{line\_\-new}}
\index{line_new@{line\_\-new}!tt_diagram.c@{tt\_\-diagram.c}}
\subsubsection[line\_\-new]{\setlength{\rightskip}{0pt plus 5cm}static \hyperlink{struct__icon__el}{icon\_\-el}$\ast$ line\_\-new (Evas $\ast$ {\em evas}, \hyperlink{struct__icon}{icon} $\ast$ {\em icon}, \hyperlink{struct__icon__el}{icon\_\-el} $\ast$ {\em el}, double {\em x1}, double {\em y1}, double {\em x2}, double {\em y2}, int {\em state})\hspace{0.3cm}{\tt  \mbox{[}static\mbox{]}}}}
\label{tt__diagram_8c_a19}


create a line connecting the file-icon with an info-item. 

\begin{Desc}
\item[Returns:]an icon-element wrapping a line \end{Desc}
\begin{Desc}
\item[Parameters:]
\begin{description}
\item[{\em evas}]- the evas to create the line-object on \item[{\em icon}]- ICON\_\-EL\_\-$\ast$() wants this \item[{\em el}]- the current icon-element \item[{\em x1}]- left boundary \item[{\em y1}]- upper boundary \item[{\em x2}]- right boundary \item[{\em y2}]- lower boundary \item[{\em state}]- state mask for the icon-elements to create\end{description}
\end{Desc}
\hypertarget{tt__diagram_8c_a24}{
\index{tt_diagram.c@{tt\_\-diagram.c}!tt_dst@{tt\_\-dst}}
\index{tt_dst@{tt\_\-dst}!tt_diagram.c@{tt\_\-diagram.c}}
\subsubsection[tt\_\-dst]{\setlength{\rightskip}{0pt plus 5cm}void tt\_\-dst (void $\ast$$\ast$ {\em data})}}
\label{tt__diagram_8c_a24}


destruct a tooltip. 

\begin{Desc}
\item[Parameters:]
\begin{description}
\item[{\em data}]- the tooltip-handle \hyperlink{tt__centre_8c_a4}{tt\_\-new()} gave us\end{description}
\end{Desc}
\hypertarget{tt__diagram_8c_a20}{
\index{tt_diagram.c@{tt\_\-diagram.c}!tt_icon_new@{tt\_\-icon\_\-new}}
\index{tt_icon_new@{tt\_\-icon\_\-new}!tt_diagram.c@{tt\_\-diagram.c}}
\subsubsection[tt\_\-icon\_\-new]{\setlength{\rightskip}{0pt plus 5cm}static \hyperlink{struct__icon}{icon}$\ast$ tt\_\-icon\_\-new (const char $\ast$ {\em label})\hspace{0.3cm}{\tt  \mbox{[}static\mbox{]}}}}
\label{tt__diagram_8c_a20}


create an \char`\"{}icon\char`\"{} (a group of info-lines) 

\begin{Desc}
\item[Returns:]an icon-struct, or NULL for out of memory \end{Desc}
\begin{Desc}
\item[Parameters:]
\begin{description}
\item[{\em label}]- the header for the group\end{description}
\end{Desc}
\hypertarget{tt__diagram_8c_a25}{
\index{tt_diagram.c@{tt\_\-diagram.c}!tt_new@{tt\_\-new}}
\index{tt_new@{tt\_\-new}!tt_diagram.c@{tt\_\-diagram.c}}
\subsubsection[tt\_\-new]{\setlength{\rightskip}{0pt plus 5cm}void$\ast$ tt\_\-new (void $\ast$ {\em data}, int {\em x}, int {\em y}, \hyperlink{struct__file__info}{file\_\-info} $\ast$ {\em fi})}}
\label{tt__diagram_8c_a25}


create a tooltip. 

no assumptions about the nature of the tooltip are made; for all the main application knows, the plugin might read the info to the user aloud.

\begin{Desc}
\item[Returns:]a handle to pass to \hyperlink{tt__centre_8c_a3}{tt\_\-dst()} when destroy the tooltip. \end{Desc}
\begin{Desc}
\item[Parameters:]
\begin{description}
\item[{\em data}]- the icon-list the tooltip is related to \item[{\em x}]- x-position of the icon the info is for \item[{\em y}]- y-position of the icon the info is for \item[{\em fi}]- the file we're displaying info for\end{description}
\end{Desc}
